%%%%%%%%%%%%%%%%%%%%%%%%%%%%%%%%%%%%%%%%%%%%%%%%%%%%%%%%%%%%%%%%%%%%%%
% Lớp 10 | Phương pháp tọa độ trong mặt phẳng (𝑂𝑥𝑦) | Đường tròn trong mặt phẳng tọa độ
% Dạng 0H9?4-3: Phương trình tiếp tuyến của đường tròn
% Số câu: 20  (sinh tự động bởi tach_cau.py; từ nay thầy sửa trực tiếp file này)
%%%%%%%%%%%%%%%%%%%%%%%%%%%%%%%%%%%%%%%%%%%%%%%%%%%%%%%%%%%%%%%%%%%%%%

%%%%%%%%%%%%%%%% TRẮC NGHIỆM 4 PHƯƠNG ÁN %%%%%%%%%%%%%%%%

% nguồn: MaHoaNH_TracNghiem_Lop10
\begin{ex}%[0H9H4-3]%[ma:be9ab330]
	Trong mặt phẳng $Oxy$, cho đường tròn $(C) \colon (x-2)^2+(y+5)^2=17$. Tiếp tuyến với $(C)$  tại điểm $A(3;-1)$ có phương trình là
	\choice 
		{\True $x+4y+1=0$} 
		 {$-x+4y+7=0$}
		{$x+4y+18=0$} 
		 {$x-6y-9=0$}
	 \loigiai{
		Tiếp tuyến với $(C)$ tại điểm $A(3;-1)$ có véc-tơ pháp tuyến là $\vec{IA}=(1;4)$ nên  có dạng $x+4y+c=0$.\\
		 Vì $A$ thuộc tiếp tuyến nên $3+4\cdot (-1)+c=0\Leftrightarrow c=1$.\\
		Vậy phương trình tiếp tuyến có dạng  $x+4y+1=0$.
	} 
\end{ex}

% nguồn: MaHoaNH_TracNghiem_Lop10
\begin{ex}%[0H9H4-3]%[ma:96583302]
	Trong mặt phẳng $Oxy$, cho đường tròn  $(C)\colon (x-3)^2+(y-1)^2=10$. Phương trình tiếp tuyến của $(C)$ tại điểm $A(4; 4)$ là
	\choice 
		{\True $x+3y-16=0$} 
		{$x+3y-4=0$} 
		 {$x-3y+5=0$}
		 {$x-3y+16=0$}
	\loigiai{ 
		Đường tròn  $(C)$ có tâm $I(3; 1)$. Điểm $A(4; 4)$ thuộc đường tròn.\\
		Tiếp tuyến của $(C)$ tại điểm $A(4; 4)$ có vectơ pháp tuyến là $\overrightarrow{IA}=(1; 3)$ nên tiếp tuyến  $d$ có phương trình dạng $x+3y+c=0$.\\
		Tiếp  tuyến $d$ đi qua $A(4; 4)$ nên $4+3\cdot 4+c=0 \Leftrightarrow c=-16$.\\
		 Vậy phương trình của $d\colon x+3y-16=0$.
	 }
\end{ex}

% nguồn: MaHoaNH_TracNghiem_Lop10
\begin{ex}%[0H9H4-3]%[ma:bb64440f]
	Trong mặt phẳng $Oxy$, cho đường tròn $(C)\colon x^2+y^2+2x-6y+5=0$. Gọi $\Delta $ là tiếp tuyến của $(C)$ tại điểm $A\left(0;1\right)$. Phương  trình tổng quát của $\Delta $ là
	\choice 
		 {$x-2y+1=0$}
		 {$x+2y-2=0$}
		{$x+2y+2=0$} 
		 {\True $x-2y+2=0$}
	\loigiai{ 
	Ta  có
	\[ 
	(C)\colon x^2+y^2+2x-6y+5=0 \Leftrightarrow (C) \colon (x+1)^2  + (y-3)^2= 5.
	 \]
	Do đó đường tròn $(C)$ có tâm $I(-1; 3)$.\\ 
	Tiếp tuyến $\Delta$ của $(C)$ đi qua  điểm $A(0;1)$, nhận $\vec{IA}= (1; -2)$ là véc-tơ pháp tuyến.\\
	Vậy phương trình tổng quát của $\Delta$ là 
	\[ 
	 1\cdot (x-0)+ (-2) \cdot (y-1)=0 \Leftrightarrow x-2y+2=0.
	\] 
	} 
\end{ex}


%%%%%%%%%%%%%%%% ĐÚNG/SAI %%%%%%%%%%%%%%%%

% nguồn: MaHoaNH_DungSai_Lop10
\begin{ex}%[0H9N4-3]%[ma:69b3b642]
	Các khẳng định sau đây, khẳng định nào đúng, khẳng định  nào sai?
	\choiceTF 
		{\True Đường thẳng  $\Delta$ là tiếp tuyến của đường tròn $(C)$ tâm $I$, bán kính $R$ khi và chỉ khi $\mathrm{d}(I,\Delta)=R$}
		{Phương trình tiếp tuyến tại điểm $M(3;4)$ với đường tròn $(C)\colon (x-1)^2+(y-2)^2=8$  là $x+y+7=0$}
		{\True Đường thẳng $d\colon x-2y+5=0$  là tiếp tuyến của đường tròn $(C_1)\colon x^2+y^2=5$}
		{Cho đường tròn $(C_2)$ tâm $J(1;-2)$ đi qua điểm $A(0;1)$. Đường thẳng $\Delta_1$ là tiếp tuyến của $(C_2)$ tại điểm $A$ có véc-tơ pháp tuyến  là $\overrightarrow{n}=(1;3)$}
	\loigiai{ 
		\begin{itemchoice}
			\itemch Đường thẳng $\Delta$ là tiếp tuyến  của đường tròn $(C)$ tâm $I$, bán kính $R$ khi và chỉ khi $\mathrm{d}(I,\Delta)=R$.
			\itemch Phương trình tiếp tuyến tại điểm  $M(3;4)$ với đường tròn $(C)$ là
			$$(x-1)(3-1)+(y-2)(4-2)=8\Leftrightarrow x+y-7=0.$$ 
			\itemch Đường tròn $(C_1)$ có tâm $O(0;0)$, bán  kính $R=\sqrt{5}$.\\
			Ta có $\mathrm{d}(O,d)=\dfrac{|0-0+5|}{\sqrt{1^2+(-2)^2}}=\sqrt{5}=R$  nên $d$ là tiếp tuyến của $(C_1)$.
			\itemch Cho đường tròn $(C_2)$ tâm $J(1;-2)$ đi qua điểm $A(0;1)$. Đường thẳng $\Delta_1$  là tiếp tuyến của $(C_2)$ tại điểm $A$ nên $\Delta_1$ vuông góc với $JA$. Suy ra $\overrightarrow{JA}=(-1;3)$ là một véc-tơ pháp tuyến của $\Delta_1$, hay $\Delta_1$ có véc-tơ pháp tuyến là $\overrightarrow{n}=(1;-3)$.
		\end{itemchoice}
	} 
\end{ex}

% nguồn: MaHoaNH_DungSai_Lop10
\begin{ex}%[0H9H4-3]%[ma:069d572a]
	Trong mặt phẳng $Oxy$, cho đường tròn $(C) \colon  x^2+y^2-4x+6y-1=0$.
	 \choiceTF
		 {Tọa độ tâm của đường tròn $(C)$ là $(-2;3)$}
		{\True Bán kính  của đường tròn $(C)$ là $R=\sqrt{14}$}
		{\True Điểm $A(5;1)$ nằm ngoài đường tròn $(C)$} 
		 {\True Phương trình tiếp tuyến của $(C)$ tại $B(4;\sqrt{10}-3)$ là $-2x-\sqrt{10}y+18-3\sqrt{10}=0$}
	\loigiai{ 
		\begin{itemchoice}
			\itemch Sai.  Vì ta có $-2a=-4 \Leftrightarrow a=2$ và $-2b=6\Leftrightarrow b=-3$. suy ra tâm $I(2;-3)$.
			\itemch  Đúng. Vì $R=\sqrt{a^2+b^2-c}=\sqrt{4+9-(-1)}=\sqrt{14}$.
			\itemch Đúng. Vì $IA=\sqrt{(5-2)^2+(1+3)^2}=5>R=\sqrt{14}$.  Suy ra $A$ nằm ngoài đường tròn.
			\itemch Đúng. Vì phương trình tiếp tuyến tại $ B$  là \[(2-4)(x-4)+(-3-\sqrt{10}+3)(y-\sqrt{10}+3)\Leftrightarrow -2x-\sqrt{10}y+18-3\sqrt{10}=0.\]
		\end{itemchoice}
	 }
\end{ex}

% nguồn: MaHoaNH_DungSai_Lop10
\begin{ex}%[0H9H4-3]%[ma:e130add1]
	Trong mặt phẳng $Oxy$, cho hai điểm $A(8;0)$, $B(0;6)$. Gọi  $I$ là trung điểm của $AB$. Xét tính đúng, sai của mỗi khẳng định sau.
	\choiceTF 
		{\True  $I(4; 3)$}
		{\True  Tam giác $OAB$ là tam giác vuông}
		{ Đường tròn ngoại tiếp tam giác  $OAB$ có bán kính $R=25$}
		{Tiếp tuyến tại $A$ với đường  tròn ngoại tiếp tam giác $OAB$ có phương trình $3x+4y-24=0$}
	 \loigiai{
		\begin{enumerate}
			\item $I(x_0; y_0)$ là trung điểm của $AB$ nên  \[\heva{&x_0=\dfrac{8+0}{2}=4\\&y_0=\dfrac{0+6}{2}=3}\Rightarrow I(4;3).\]
			\item Vì  $A(8;0)\in Ox$, $B(0;6)\in Oy$ nên tam giác $OAB$ vuông tại $O(0;0)$.
			\item Bán kính của đường tròn đi qua $O$, $A$, $B$ là  $R=\dfrac{1}{2}AB=\dfrac{1}{2}\sqrt{(0-8)^2+(6-0)^2}=5$.
			\item Ta  có $\vec{IA}(4; -3)$.\\
			Tiếp tuyến tại $A$ với đường tròn ngoại tiếp tam giác $OAB$ đi qua $A$ và nhận $\vec{IA}$ làm véc-tơ pháp  tuyến nên có phương trình
			\[4(x-8)-3y=0\Leftrightarrow  4x-3y-32=0.\]
		\end{enumerate}
	 }
\end{ex}

% nguồn: MaHoaNH_DungSai_Lop10
\begin{ex}%[0H9H4-3]%[ma:f3f5457e]
	Các khẳng định sau đây, khẳng định nào đúng, khẳng định nào sai? 
	\choiceTF 
		{Đường thẳng $\Delta\colon x+my-2=0$ là tiếp tuyến của đường  tròn $(C)\colon (x-1)^2+(y+2)^2=4$ khi và chỉ khi $m=-\dfrac{5}{2}$ hoặc $m=\dfrac{1}{2}$}
		{\True Số đường thẳng đi qua điểm $M(5;6)$ và tiếp xúc với đường tròn $(C)\colon (x-1)^2+y^2=4$ là $2$} 
		{\True Phương trình  tiếp tuyến của đường tròn $(C)\colon x^2+y^2-2x-4y+3=0$ tại giao điểm của $(C)$ với trục tung là $x+y-1=0$ hoặc $x-y+3=0$}
		{\True Đường thẳng vuông góc với đường thẳng  $d\colon x+y-1=0$ và tiếp xúc với đường tròn $(C)\colon x^2+(y-3)^2=2$ có phương trình $x-y+1=0$ hoặc $x-y+5=0$}
	 \loigiai{
		\begin{itemchoice}
			\itemch Đường tròn $(C)\colon (x-1)^2+(y+2)^2=4$ có  tâm $I(1;-2)$, bán kính $R=2$.\\
			Đường thẳng $\Delta\colon x+my-2=0$ là tiếp tuyến của đường tròn  $(C)$ khi và chỉ khi
			$$\mathrm{d}(I,\Delta)=R  \Leftrightarrow \dfrac{|1-2m-2|}{\sqrt{1+m^2}}=2 \Leftrightarrow 4m^2+4m+1=4m^2+4 \Leftrightarrow m=\dfrac{3}{4}.$$
			\itemch Đường tròn $(C)$ có tâm  $I(1;0)$, bán kính $R=2$.\\
			 Ta có $\overrightarrow{IM}=(4;6) \Rightarrow IM=\sqrt{4^2+6^2}=2\sqrt{13}>2=R$ nên điểm $M$ nằm ngoài đường tròn $(C)$.\\
			Vậy qua $M$ có $2$ đường thẳng  là tiếp tuyến của $(C)$.
			\itemch Đường tròn  $(C)\colon x^2+y^2-2x-4y+3=0$ có tâm $I(1;2)$ và bán kính $R=\sqrt{1^2+2^2-3}=\sqrt{2}$.\\
			Gọi $A(0;a)$ thuộc trục $Oy$  là giao điểm của $(C)$ với trục tung. Khi đó ta có $A\in (C)$ nên
			$$a^2-4a+3=0 \Leftrightarrow \hoac{&a=1\\&a=3}  \Rightarrow \hoac{&A(0;1)\\&A(0;3).}$$
			\begin{itemize}
				\item Với  điểm $A(0;1)$.\\
				Tiếp tuyến của $(C)$ tại $A(0;1)$  có phương trình
				 $$-1(x-0)-1(y-1)=0\Leftrightarrow x+y-1=0.$$
				\item  Với điểm $A(0;3)$.\\
				Tiếp tuyến của $(C)$ tại $A(0;3)$ có phương  trình
				$$-1(x-0)+1(y-3)=0\Leftrightarrow  -x+y-3=0.$$
			\end{itemize}
			 \itemch Đường tròn $(C)\colon x^2+(y-3)^2=2$ có tâm $I(0;3)$, bán kính $R=\sqrt{2}$.\\
			Đường thẳng $d$ vuông  góc với đường thẳng $x+y-1=0$ nên có phương trình dạng $x-y+c=0$.\\
			Đường  thẳng $d$ là tiếp tuyến của $(C)$ khi và chỉ khi
			$$\mathrm{d}(I,d)=R \Leftrightarrow \dfrac{|0-3+c|}{\sqrt{1^2+(-1)^2}}=\sqrt{2} \Leftrightarrow  |c-3|=2 \Leftrightarrow \hoac{&c=1\\&c=5.}$$
			Vậy phương trình đường thẳng $d$ là  $x-y+1=0$ hoặc $x-y+5=0$.
		\end{itemchoice}
	 }
\end{ex}

% nguồn: MaHoaNH_DungSai_Lop10
\begin{ex}%[0H9V4-3]%[ma:d8a44106]
	Trong mặt phẳng $Oxy$, cho đường tròn $(C)$ có phương trình $x^2+y^2-6x+2y+6=0$  và hai điểm $A(1;-1)$, $B(1;3)$.
	 \choiceTF
		 {\True Điểm $A$ thuộc đường tròn}
		{Điểm $B$ nằm trong  đường tròn}
		 {\True $x=1$ phương trình tiếp tuyến của $(C)$ tại điểm $A$}
		{Qua $B$ kẻ được hai  tiếp tuyến với $(C)$ có phương trình là $x=1$; $3x+4y-12=0$}
	 \loigiai{
		 Đường tròn $(C)$ có tâm $I(3;-1)$ bán kính $R=\sqrt{9+1-6}=2$.
		\begin{itemchoice}
			\itemch  Ta có $IA=2=R$ nên $A\in (C)$.
			 \itemch Ta có $IB=2\sqrt{5} > R$ suy ra điểm $B$ nằm ngoài đường tròn.
			 \itemch Tiếp tuyến của $(C)$ tại điểm $A$ nhận $\overrightarrow{AI}=(2;0)$ làm véc-tơ pháp tuyến nên có phương trình là $2(x-1)+0(y+1)=0$ hay $x=1$.
			\itemch Phương trình đường thẳng $\Delta$ đi qua $B$ có dạng  $a(x-1)+b(y-3)=0$ (với $a^2+b^2 \ne 0$) hay $ax+by-a-3b=0$.\\
			Đường thẳng $\Delta$ là tiếp tuyến của đường tròn khi  và chỉ khi
			\[ 
			\mathrm{d}(I,\Delta)=R \Leftrightarrow \dfrac{|3a-b-a-3b|}{\sqrt{a^2+b^2}}=2\Leftrightarrow (a-2b)^2=a^2+b^2\Leftrightarrow 3b^2-4ab=0\Leftrightarrow \hoac{& b=0 \\ &  3b=4a.}
			 \]
			Với $b=0$, chọn $a=1$; phương trình tiếp tuyến  là $x=1$.\\
			Với $3b=4a$, chọn $a=3\Rightarrow b=4$; phương trình tiếp  tuyến là $3x+4y-15=0$.\\
			Vậy qua $B$ kẻ được hai tiếp tuyến với $(C)$ có phương trình  là $x=1$; $3x+4y-15=0$.\\
		\end{itemchoice}
	} 
\end{ex}


%%%%%%%%%%%%%%%% TRẢ LỜI NGẮN %%%%%%%%%%%%%%%%

% nguồn: MaHoaNH_TraLoiNgan_Lop10
\begin{bt}%[0H9H4-3]%[ma:c44fed4d]
	Trong mặt phẳng $Oxy$, cho đường tròn $(C)$ có phương trình $(x-3)^2+(y-5)^2=100$ và tiếp xúc với đường thẳng $\Delta$ có phương trình $4x+3y+m=0$ với $m$ là tham  số dương. Tìm giá trị của $m$.
	 \shortans{23}
	\loigiai{ 
		Phương trình đường  tròn $(C)$ có tâm $I(3;5)$, bán kính $R=10$.\\
		$$\mathrm{d}(I,\Delta)=R\Leftrightarrow\dfrac{|4\cdot3+3\cdot5+m|}{\sqrt{4^2+3^2}}=10\Leftrightarrow |27+m|=50\Leftrightarrow\hoac{&27+m=50\\&27+m=-50}\Leftrightarrow\hoac{&m=23\\&m=-77}.$$ 
		 Vậy $m=23$.
	} 
\end{bt}

% nguồn: MaHoaNH_TraLoiNgan_Lop10
\begin{bt}%[0H9H4-3]%[ma:cda67c51]
	 Trong hệ trục tọa độ $Oxy$, cho đường tròn $(C)\colon (x-1)^2+(y+2)^2=8$. Phương trình tiếp tuyến $d$ của $(C)$ tại điểm $A(3;-4)$ có dạng $x+by+c=0$ với $b, c\in\mathbb{R}$. Khi đó $b^2+c^2$ bằng
	\shortans{50} 
	 \loigiai{
		 Đường tròn $(C)$ có tâm $I(1;-2)$.\\
		Tiếp tuyến của đường tròn $(C)$ tại $A$ là đường thẳng đi qua $A$ và nhận $\overrightarrow{IA}=(2;-2)$ làm véc-tơ pháp tuyến.  Phương trình tiếp tuyến là $$2(x-3)-2(y+4)=0\Leftrightarrow x-y-7=0.$$
		Vậy  $b=-1, c=-7$ nên $b^2+c^2=50$.
	 }
\end{bt}

% nguồn: MaHoaNH_TraLoiNgan_Lop10
\begin{bt}%[0H9H4-3]%[ma:a0d5537c]
	Trong mặt phẳng  $Oxy$, gọi $\Delta$ là tiếp tuyến với đường tròn $x^2+y^2=5$ tại điểm $M(1;2)$. Hệ số góc của $\Delta$ bằng bao nhiêu? \shortans{-0,5}
	\loigiai{Ta có $M(1;2)$ thuộc  đường tròn.\\
		 Đường tròn đã cho có tâm $O(0;0)$.\\
		Tiếp tuyến với đường tròn tại $M(1;2)$, có véc-tơ pháp tuyến $\overrightarrow{OM}=(1;2)$ nên  có phương trình
		\[1(x-1)+2(y-2)=0\Leftrightarrow  x+2y-5=0.\]
		Vậy tiếp tuyến cần tìm có hệ số  góc $k=-\dfrac{1}{2}=-0{,}5$.
	 }
\end{bt}

% nguồn: MaHoaNH_TraLoiNgan_Lop10
\begin{bt}%[0H9H4-3]%[ma:9b86d038]
	 Trong mặt phẳng $Oxy$, cho đường tròn $(C)\colon x^2+y^2+2x-4y-4=0$. Biết rằng qua điểm $A(2;m)$ chỉ có một tiếp tuyến với $(C)$. Giá trị của $m$ bằng bao nhiêu?
	\shortans{2} 
	 \loigiai{
		Qua điểm $A$ chỉ có một tiếp tuyến với  đường tròn
		 \[
		(C)\Leftrightarrow A\in  (C) \Leftrightarrow 2^2+m^2+2\cdot 2-4m-4=0\Leftrightarrow m^2-4m+4=0\Leftrightarrow (m-2)^2=0\Leftrightarrow m=2.
		\] 
	 }
\end{bt}

% nguồn: MaHoaNH_TraLoiNgan_Lop10
\begin{bt}%[0H9H4-3]%[ma:d3dd1ac1]
	Trong mặt phẳng $Oxy$, cho đường tròn $(C)\colon (x-1)^2+(y+2)^2=8$. Đường thẳng $d$ có phương trình dạng $ax+by-7=0$ là tiếp tuyến của  $(C)$ tại điểm $M(3;-4)$. Khi đó $a-b$ bằng bao nhiêu?
	\shortans{2} 
	\loigiai{ 
		Đường tròn  $(C)$ có tâm $I(1;-2)$.\\
		Phương trình tiếp tuyến của $(C)$ tại  điểm $M(3;-4)$ là
		 $$(3-1)(x-3)+(-4+2)(y+4)=0\Leftrightarrow2x-2y-14=0\Leftrightarrow x-y-7=0.$$
		Vậy $a=1$, $b=-1$  và $a-b=2$.
	 }
\end{bt}

% nguồn: MaHoaNH_TraLoiNgan_Lop10
\begin{bt}%[0H9H4-3]%[ma:7adfc4f6]
	Trong mặt phẳng $Oxy$, cho đường tròn $(C)\colon (x-7)^2+(y+3)^2=100$. Phương trình  tiếp tuyến của $(C)$ tại điểm $A(1;5)$ có dạng $ax+by+17=0$. Giá trị của $a^3-b^2$ là
	 \shortans{11}
	 \loigiai{
		\begin{itemize}
			\item Đường tròn $(C)\colon (x-7)^2+(y+3)^2=100$ có  tâm $I(7;-3)$.
			\item Phương trình tiếp tuyến của $(C)$ tại điểm $A(1;5)$  là
			$$(7-1)(x-1)+(-3-5)(y-5)=0\Leftrightarrow3x-4y+17=0.$$ 
			Vậy  $a^3-b^2=3^3-(-4)^2=11$.
		\end{itemize}
	 }
\end{bt}

% nguồn: MaHoaNH_TraLoiNgan_Lop10
\begin{bt}%[0H9V4-3]%[ma:4256b19c]
	Trong mặt phẳng tọa độ $Oxy$, cho đường thẳng $d \colon mx-y+2m-1=0$ (với $m$ là tham số) và đường tròn $(C) \colon (x-3)^2+(y+1)^2=20$. Có bao nhiêu giá trị dương của tham số $m$ để $d$ tiếp xúc với  $(C)$?
	 \shortans {$1$}
	 \loigiai{
		Đường tròn $(C)$ có  tâm $I(3;-1)$, bán kính $R=2\sqrt{5}$.\\
		Đường thẳng $d$ tiếp xúc với đường tròn  $(C)$ khi và chỉ khi
		$$\mathrm{d}(I;d)=R \Leftrightarrow  \dfrac{|3m+1+2m-1|}{\sqrt{m^2+1}}=2\sqrt{5}
		\Leftrightarrow |5m|=\sqrt{20(m^2+1)} \Leftrightarrow  m^2=4 \Leftrightarrow m=\pm 2.$$
		Vì $m>0$ nên  $m=2$.
	} 
\end{bt}

% nguồn: MaHoaNH_TraLoiNgan_Lop10
\begin{bt}%[0H9V4-3]%[ma:4797739b]
	Trong mặt phẳng $Oxy$, cho đường tròn $(C)\colon x^2+y^2-4x+4y-1=0$ và đường thẳng $d \colon 2x+3y+4=0$. Biết rằng có hai đường thẳng $\Delta$ vuông góc với $d$ đồng thời là tiếp tuyến của $(C)$. Biết hai tiếp tuyến này cắt trục tung lần lượt tại hai điểm $M$, $N$. Tổng tung độ hai điểm $M$, $N$  bằng bao nhiêu?
	\shortans{-5} 
	 \loigiai{
		Đường tròn $(C)$  có tâm $I(2;-2)$, bán kính $R=3$.\\
		Vì $\Delta \perp  d$ nên phương trình $\Delta$ có dạng $3x-2y+c=0$.\\
		Đường thẳng $\Delta$ tiếp  xúc đường tròn $(C)$ khi và chỉ khi $d(I,\Delta)=3$
		\[\Leftrightarrow \dfrac{|3.2-2(-2)+c|}{\sqrt{9+4}}=3\Leftrightarrow|c+10|=3\sqrt{13}  \Leftrightarrow c=-10\pm 3\sqrt{13}.
		 \]
		Hai tiếp tuyến thỏa mãn là  $\Delta \colon 3x-2y-10\pm 3\sqrt{13}=0$.\\
		Hai đường thẳng này cắt trục tung tại $M\left(0; \dfrac{-5 + 3\sqrt{13}}{2}\right)$ và  $N\left(0; \dfrac{-5 - 3\sqrt{13}}{2}\right)$.\\
		Tổng tung độ hai điểm trên bằng  $-5$.
	} 
\end{bt}

% nguồn: MaHoaNH_TraLoiNgan_Lop10
\begin{bt}%[0H9V4-3]%[ma:1be0bb3b]
	Trong mặt phẳng $Oxy$, cho điểm $M(2;-2)$ và đường tròn $(C)\colon x^2+y^2=4$. Gọi $T_1$, $T_2$ lần lượt là tiếp điểm của các tiếp tuyến kẻ từ $M$  đến $(C)$. Nếu phương trình $T_1T_2$ là $3x+by+c=0$ thì giá trị biểu thức $4b+c$ bằng
	\shortans{-18} 
	\loigiai{ 
		Đường tròn $(C)$ có tâm là  $O$.\\
		Vì $T_1$, $T_2$ là các tiếp điểm của tiếp tuyến kẻ từ $M$ đến  $(C)$ nên $MT_1\perp OT_1, MT_1\perp OT_1$.\\
		Vậy $T_1$, $T_2$  thuộc đường tròn $(C')$ có đường kính là $MO$.\\
		Gọi $I$ là trung điểm $MO$ thì $I(1;-1)$, bán  kính $R'=\dfrac{1}{2}MO=\sqrt{2}$.\\
		Phương trình đường tròn $(C')$ là $(x-1)^2+(y+1)^2=2$.\\ 
		Tọa độ hai điểm $T_1$, $T_2$ là nghiệm  của hệ
		$$\heva{& x^2+y^2=4 \\ & (x-1)^2+(y+1)^2=2}\Leftrightarrow \heva{& x^2+y^2=4 \\ & x^2+y^2-2x+2y=0}\Rightarrow -2x+2y+4=0\Leftrightarrow  3x-3y-6=0.$$
		Vậy phương trình đường thẳng $T_1T_2$  là $3x-3y-6=0$.\\
		Vậy $b=-3, c=-6$ nên $4b+c=-18$. 
	 }
\end{bt}

% nguồn: MaHoaNH_TraLoiNgan_Lop10
\begin{bt}%[0H9V4-3]%[ma:4ee022be]
	Trong mặt phẳng $Oxy$, cho hai điểm $A(4;3)$ và $B(2;1)$ cùng thuộc đường tròn $(C)$. Các tiếp tuyến của $(C)$ tại $A$ và $B$  cắt nhau tại điểm $I$ thuộc trục hoành. Tìm hoành độ điểm $I$.
	 \shortans{5}
	 \loigiai{
		\begin{itemize}
			\item Gọi $I(x;y)$. Vì $I$ thuộc trục  hoành nên $I(x;0)$.
			\item Vì các tiếp tuyến của $(C)$ tại $A$  và $B$ cắt nhau tại điểm $I$ nên $AI=BI$.
			 \item Ta có $AI=\sqrt{(x-4)^2+9}$ và $BI=\sqrt{(x-2)^2+1}$.
			 \item Giải phương trình
			\allowdisplaybreaks 
			\begin{eqnarray*}
				&&\sqrt{(x-4)^2+9}=\sqrt{(x-2)^2+1}\\ 
				&\Rightarrow&  x^2-8x+25=x^2-4x+5\\
				&\Leftrightarrow&4x=20\\ 
				&\Leftrightarrow&x=5. 
			\end{eqnarray*}
			 Thay giá trị $x=5$ vào phương trình thỏa mãn.
			\item Vậy tọa độ  điểm $I(5;0)$.
		\end{itemize}
	} 
\end{bt}

% nguồn: MaHoaNH_TraLoiNgan_Lop10
\begin{bt}%[0H9C4-3]%[ma:7f76fb19]
	Trong mặt phẳng tọa độ $Oxy$, cho đường tròn $(C)\colon (x+1)^2+(y-2)^2 = 5$ và đường thẳng $\Delta\colon x-y+1=0$. Gọi $M(a;b)$ với  $a>0$ là điểm trên $\Delta$ sao cho từ $M$ có thể kẻ được hai tiếp tuyến đến $(C)$ có tiếp điểm là $A$ và $B$ thỏa mãn $\widehat{AMB} = 60^\circ$. Khi đó tổng $20a^3 + 9b^2$ bằng
	 \shortans{684}
	\loigiai{ 
	 \immini{
		Đường tròn $(C)$ có tâm $I(-1;2)$ và bán kính  $R =\sqrt{5}$.\\
		Giả sử $M(a;a+1)\in\Delta\colon x-y+1=0$. \\ 
		Xét $\triangle IAM$ vuông tại $A$, có $\widehat{AMI}=\dfrac{1}{2}\widehat{AMB}  = 30^\circ$.\\
		Suy ra $IM =\dfrac{IA}{\sin\widehat{AIM}} = \dfrac{R}{\sin  30^\circ}=2\sqrt{5}$.\\
		Vậy ta có phương  trình
		 $$(a+1)^2 +(a-1)^2 = \left(2\sqrt{5}\right)^2 \Leftrightarrow \hoac{&a=3\\&a=-3\quad(\text{loại vì } a>0).}$$
		 }{
		\begin{tikzpicture}[scale=.6,font=\footnotesize,line join=round,line cap=round,>=stealth]
				\path 
					(5,0) coordinate (I) 
					(-3,0) coordinate  (M)
					(3.875,2.781,0) coordinate  (A)
					(3.875,-2.781)  coordinate (B);
				\draw[] (I) circle (3cm); 
				\draw[]  (I)--(M)--(A)--(I)--(B)--(M)
					(A)--(B); 
				\foreach \x/\g in{A/100,B/-100,I/0,M/-180} 
				 \fill[black](\x) circle (1pt)
					($(\x)+(\g:4mm)$)  node{$\x$};
				\draw pic[draw,angle  radius=2mm] {right angle = M--A--I};
				\draw pic[draw,angle radius=2mm]  {right angle = M--B--I};
			\end{tikzpicture}
		 }
		\noindent  Với $a= 3$ thì $M(3;4)$. Suy ra $a = 3$, $b = 4$ nên $20a^3 + 9b^2 = 684$.
	 }
\end{bt}


%%%%%%%%%%%%%%%% TỰ LUẬN %%%%%%%%%%%%%%%%

% nguồn: MaHoaNH_TuLuan_Lop10
\begin{bt}%[0H9V4-3]%[ma:2614124a]
	Trong mặt phẳng $Oxy$, cho đường tròn $(C)\colon (x-6)^2+(y-7)^2=25$. Gọi $d$ là đường thẳng tiếp xúc với đường tròn $(C)$ và song song với đường thẳng $\Delta  \colon x+2y-5\sqrt{5}-20=0$. Đường thẳng $d$ cắt đường thẳng $x + 5\sqrt{5} = 0$ tại điểm có tung độ bằng bao nhiêu?
	 \loigiai{
		Đường tròn $(C)$ có tâm $I(6;7)$, bán  kính $R=5$.\\
		Đường thẳng  $d$ song song với $\Delta$ nên có dạng: $x+2y+m=0, (m\ne -5\sqrt{5}-20)$.\\
		Mặt khác, đường thẳng $d$  tiếp xúc với đường tròn $(C)$ nên\\
		\[ 
		\mathrm{d}(I,\Delta)=\dfrac{|6+2\cdot 7+m|}{\sqrt{1^2+2^2}}=\dfrac{|20+m|}{\sqrt{5}}=5. 
		 \]
		Giải ra  được $m=5\sqrt{5}-20$ (thỏa mãn) hoặc $m=-5\sqrt{5}-20$ (loại).\\
		 Suy ra $d\colon x+2y + 5\sqrt{5}-20=0$.\\
		Do  đó đường thẳng $d$ cắt đường thẳng $x + 5\sqrt{5} = 0$ tại điểm có tung độ thỏa mãn $2y - 20 = 0 \Leftrightarrow y = 10$.
	 }
\end{bt}
